% ============================================================
% Columbia Climate School Postdoctoral Research Program -- Research Proposal
% Format (program rules): max 5 pages excl. references; double-spaced;
% 1-inch margins; minimum 11pt standard font; title required.
% ============================================================
\documentclass[11pt]{article}
\usepackage[margin=1in]{geometry}
\usepackage{setspace}
\usepackage{amsmath,amssymb}
\usepackage{booktabs}
\usepackage[T1]{fontenc}
\usepackage{lmodern}
\usepackage[hidelinks]{hyperref}
\usepackage{natbib}

\begin{document}
\doublespacing

\begin{center}
{\large\bfseries Riders on the Storm}\\[2pt]
{\itshape Platform Market Power and the Incidence of Storms in New York City's Ride-Hail Market, 2019--2026}
\end{center}

\noindent\textbf{Prepared by:} Emil Kee-Tui [affiliation] \textperiodcentered{} \textbf{Co-author:} Xibo Wan [affiliation]\\
\textbf{Proposed mentor:} Dr.\ Jacqueline M.\ Klopp, Director, Center for Sustainable Urban Development, Columbia Climate School; [second mentor, to be confirmed]

\section{Motivation and Research Question}

Hurricanes and tropical storms are increasing in frequency due to climate change and are a major cost associated with the damages of climate. Hurricanes, flood subways and endanger drivers, bikers, and pedestrians. But rideshare platforms and drivers that continue to offer rides can surge prices during these conditions since they face less competition from other modes of transport and drivers and more demand for evacuation. During Hurricane Ida (1 September 2021), fares reached about three times the TLC pay floor. Despite the drivers undertaking all the risk of driving during a hurricane the platforms share of fare increased: between 9 pm, when the heaviest rain began, and midnight the driver's share of fare fell from 91\% to 73\% at Uber and from 72\% to 45\% at Lyft, against a normal September level of 85\% at Uber and 73\% at Lyft.

We study how hurricanes and tropical storms affect ridesharing fares, rides taken, surge pricing, and the distribution of welfare among riders, drivers, and platforms in the New York City ridesharing market from 2019 to 2026. Understanding the impact of hurricanes and tropical storms on the NYC rideshare market is important as extreme weather events are projected to increase in intensity due to climate change and the market facilitates large amounts of transactions and economic value, in 2026 it was valued at \$5.2 billion, employed over 80,000 drivers, and facilitated hundreds of thousands of daily trips.

\section{Related Literature and Theoretical Framing}

Surge pricing reduces wasteful driver search \citep{castillo2017} and raises rider but lowers driver surplus \citep{castillo2025}; rain shifts taxi demand and labor supply \citep{farber2015,brodeur2018}; platforms hold wage-setting power \citep{dube2020}; and hurricane damage is large \citep{kahn2005,hsiang2014,deryugina2017}. But these literatures study routine, reversible imbalances, where riders have time and alternatives, or aggregate, longer-run costs. None study surge pricing during extreme weather when demand is acutely inelastic (life-safety and evacuation trips), competition from other forms of transport and other drivers and driver's outside work options are simultaneoulsy affected, and anti-price-gouging law may be legally activated.

We analyze storms through platform market power in the rider market and the driver market. When both markets are competitive, platforms offer drivers a share of fares equal to their opportunity cost of driving. Market power in the rider market allows platforms to charge a higher fare; wage-setting power over drivers allows them to reduce the driver's share of fare and capture a greater percentage of fare increases. With demand elasticity $\varepsilon_D$, driver supply elasticity $\varepsilon_S$ and platform cost $c$, the driver's share is
\begin{equation}
\text{DriverShare} = \frac{\text{driver\_pay}}{\text{base\_passenger\_fare}}.
\end{equation}

Hurricanes increase platform power in the rider market if they reduce competition from alternative transport (subway closures, walking becoming more dangerous than driving) or if rideshares are especially suited for evacuations. They reduce drivers' market power if drivers' outside employment options close, and increase it if drivers demand higher wages to accept the risk of driving during a hurricane. 

During a hurricane the pass-through elasticity of the driver's pay and the surge is

\begin{equation}
\frac{\partial \text{driver\_pay\_share}}{\partial \text{surge}}.
\end{equation}

\section{Data and Variables}

The data discussed in this section has been accessed and downloaded. Table~\ref{tab:vars} defines the variables used in the analysis. The NYC High Volume For-Hire Vehicle (HVFHV) trip records, which rideshare companies must legally file with the TLC, report the request, pickup and dropoff times, taxi zones, distance, duration, fare paid, and driver pay for every Uber, Lyft, and Via trip from 2019 to 2026. Lacking a surge multiplier, we measure surge as fare per mile relative to the zone's non-storm baseline for that hour of the week, and aggregate trips to taxi zone $\times$ hour ($z,t$) and origin--destination pair $\times$ hour. Yellow and green taxi records give the taxi share; the timing of each storm $s$ is measured with NOAA HURDAT2; storm surge (SLOSH, a bathtub elevation model and tide gauges), FEMA floodplains and 311 complaints show where water reached; MTA alerts record each bus and subway delay or closure; and cellphone data measure population present and net outflow by zone and hour.

\begin{table}[h]
\centering
\begin{singlespace}
\footnotesize
\renewcommand{\arraystretch}{0.9}
\setlength{\abovecaptionskip}{4pt}
\begin{tabular}{ll}
\toprule
\multicolumn{2}{l}{\textit{Outcomes} ($Y_{z,s,t}$)} \\
\quad SurgeProxy & Fare per mile / zone's non-storm hour-of-week baseline (HVFHV) \\
\quad DriverShare, TakeRate & driver\_pay / base\_passenger\_fare; $1 -$ DriverShare (HVFHV) \\
\quad WaitTime (SupplyProxy) & Pickup time $-$ request time; dispatch delay (HVFHV) \\
\quad TaxiShare & Taxi trips / (taxi $+$ HVFHV trips) (TLC taxi records) \\
\midrule
\multicolumn{2}{l}{\textit{Treatment and timing}} \\
\quad FloodExposed$_{z,s}$ & Continuous: rainfall, surge depth, floodplain, 311 reports \\
\quad EventHour$_s$ & Onset of peak rainfall over NYC (HURDAT2, radar) \\
\quad $X_{z,t}$ & Hourly radar rainfall in the zone \\
\midrule
\multicolumn{2}{l}{\textit{Mechanisms}} \\
\quad DemandProxy & Population present and net outflow vs.\ pre-storm baseline (cellphone) \\
\quad TransitDisruption & Subway and bus delays and closures in the zone (MTA alerts) \\
\midrule
\multicolumn{2}{l}{\textit{Instrument and controls}} \\
\quad TrackDistance, TrackAngle & Zone distance to storm track; track angle to coastline \\
\quad Elevation, CoastDistance & Zone elevation; distance to coastline \\
\bottomrule
\end{tabular}
\end{singlespace}
\caption{Variables used in the empirical strategy.}
\label{tab:vars}
\end{table}

\section{Empirical Strategy}

\textbf{Multi-storm stacked event study.} We pool every tropical storm and hurricane between 2019--2026:
\begin{equation}
Y_{z,s,t} = \alpha_{z,s} + \gamma_{s,t} + \sum_{k \neq -1} \beta_k \cdot \text{FloodExposed}_{z,s} \times \mathbb{1}[t = \text{EventHour}_s + k] + X_{z,t}\,\delta + \varepsilon_{z,s,t}
\end{equation}
The unit is taxi zone $\times$ hour, stacked by storm. $\beta_k$ traces out the dynamic effect $k$ hours from the event hour, and leads capture anticipatory evacuation demand. Outcomes are those listed in Table~\ref{tab:vars}. Stacking avoids staggered-adoption bias in two-way fixed effects \citep{cengiz2019,callaway2021} and tests whether pricing algorithms changed across storms.

\textbf{Supply--demand decomposition.} $\text{SurgeProxy} = \beta_1 \cdot \text{DemandProxy} + \beta_2 \cdot \text{SupplyProxy} + \text{FE}$, with the proxies in Table~\ref{tab:vars}. Second-stage test: does surge at $t$ actually raise supply at $t+1$? Within origin--destination pairs, $\log \text{share} = \log \text{pay} - \log \text{fare}$ splits each share effect into fare and pay.

\textbf{Storm-track instrument.} Flooded zones may differ in supply, income or drainage even without a storm, so we instrument flood exposure with the zone's distance from the storm track and the track's angle to the coastline, controlling for elevation, distance to coast and rainfall.

\textbf{Inference and welfare.} Standard errors are clustered by zone, with storm-level wild cluster bootstrap \citep{cameron2008}; storms with and without a governor's emergency declaration (General Business Law \S396-r) are compared. Welfare separates transfers ($\Delta$Fare $\times$ TripVolume, split into $\Delta$DriverPay and $\Delta$PlatformMargin), deadweight loss from unmet demand (a taxi vs.\ HVFHS vs.\ no-trip choice model) and driver risk compensation.

\section{Preliminary Evidence}

\textbf{Sample.} All Uber, Lyft and Via trips from January to December 2021, excluding Access-a-Ride trips, trips with a fare or driver pay of zero or less, fares above \$1,000, distances of zero or above 100 miles, durations under one minute or over six hours, and unknown pickup or dropoff zones: 171.4 million trips. Trips are collapsed to 106.4 million cells (platform $\times$ pickup zone $\times$ dropoff zone $\times$ pickup hour $\times$ dropoff hour) and weighted by trips per cell, so the estimates equal trip-level OLS. Rainfall is radar rainfall averaged over each taxi zone and hour.

\textbf{Specification.} For trip $i$ on platform $p$, from origin zone $o$ in hour $h$ to destination zone $d$ in hour $h'$,
\begin{multline}
\text{DriverShare}_i = \sum_{b} \beta_b \mathbb{1}[R_{o,h} \in b] + \sum_{b} \theta_b \mathbb{1}[R^{6}_{o,h} \in b] + \sum_{b} \lambda_b \mathbb{1}[R_{d,h'} \in b] + \sum_{b} \mu_b \mathbb{1}[R^{6}_{d,h'} \in b] \\
+ \alpha_p + \phi_{o,\text{hod}(h)} + \psi_{d,\text{hod}(h')} + \tau_{\text{month}} + \omega_{\text{dow}} + \varepsilon_i
\end{multline}
where $R$ is rain in the current hour (11 bins from 0.1 to over 40 mm/h) and $R^{6}$ is total rain over the previous six hours (bins up to over 75 mm), each relative to dry ($<0.1$ mm). The fixed effects are platform, origin zone $\times$ hour of day, destination zone $\times$ hour of day, month and day of week; standard errors are clustered by pickup zone and date. The same equation with log fare and log pay as outcomes splits each share effect into fare and pay.

\textbf{Results.} Fare and driver pay move one for one in ordinary rain. After more than 50 mm of rain at the origin over the previous six hours, fare outpaces pay and the driver share falls 7--8 percentage points ($\theta = -0.083$ for 50--75 mm and $-0.074$ above 75 mm, both $p<0.01$). The TLC pay floor bound on about 15\% of Uber and 62\% of Lyft trips in September 2021.

\section{Qualifications}

[Applicant: PhD field, institution, year.] [Xibo Wan: background.] We have built the full pipeline in Python and an exact batched Frisch--Waugh--Lovell estimator in R that matches \texttt{fixest} to $10^{-9}$ on 106 million observations.

\section{Timeline and Deliverables}

\begin{enumerate}
\item Months 1--4: finalize the storm sample and pull matching TLC/HVFHS months; build the flood exposure layer (SLOSH/FEMA, tide gauge and elevation); construct the zone-hour panel.
\item Months 5--12: estimate the event study/DiD, decomposition and pass-through elasticity; Paper 1 submitted; Fall symposium.
\item Months 13--24 (extension): storm-track IV and welfare model; Paper 2 on who bears storm costs; policy brief for the TLC and MTA.
\end{enumerate}

\section{Expected Contributions}

\begin{itemize}
\item The first test of surge-pricing efficiency under acutely inelastic, life-safety demand combined with physically constrained supply.
\item A structural pass-through elasticity (driver pay share vs.\ surge) as the headline estimate, not just a reduced-form price effect, with identification beyond a flood DiD: a storm-track IV and a multi-storm stacked design.
\item For adaptation and sustainable development: evidence on whether algorithmic pricing helps or hinders mobility in extreme weather and who bears the cost, transferable to flood-exposed cities of the Global South wherever trip reporting is mandated, and relevant to TLC pay rules, price-gouging law and transit-suspension decisions.
\end{itemize}

\section{Mentoring and Requested Support}

Dr.\ Klopp (CSUD) on transport governance and agency partners for the policy brief; a second mentor in economics [to be confirmed]; feedback on the framework and stacked design before estimation at scale.

\clearpage
\begin{singlespace}
\begin{thebibliography}{99}
\bibitem[Brodeur and Nield(2018)]{brodeur2018} Brodeur, A. and K. Nield (2018). ``An Empirical Analysis of Taxi, Lyft and Uber Rides: Evidence from Weather Shocks in NYC.'' \textit{Journal of Economic Behavior \& Organization} 152: 1--16.
\bibitem[Callaway and Sant'Anna(2021)]{callaway2021} Callaway, B. and P. H. C. Sant'Anna (2021). ``Difference-in-Differences with Multiple Time Periods.'' \textit{Journal of Econometrics} 225(2): 200--230.
\bibitem[Cameron et~al.(2008)]{cameron2008} Cameron, A. C., J. B. Gelbach and D. L. Miller (2008). ``Bootstrap-Based Improvements for Inference with Clustered Errors.'' \textit{Review of Economics and Statistics} 90(3): 414--427.
\bibitem[Castillo(2025)]{castillo2025} Castillo, J. C. (2025). ``Who Benefits from Surge Pricing?'' \textit{Econometrica} 93(5): 1811--1854.
\bibitem[Castillo et~al.(2017)]{castillo2017} Castillo, J. C., D. Knoepfle and E. G. Weyl (2017). ``Surge Pricing Solves the Wild Goose Chase.'' \textit{Proceedings of the 2017 ACM Conference on Economics and Computation}: 241--242.
\bibitem[Cengiz et~al.(2019)]{cengiz2019} Cengiz, D., A. Dube, A. Lindner and B. Zipperer (2019). ``The Effect of Minimum Wages on Low-Wage Jobs.'' \textit{Quarterly Journal of Economics} 134(3): 1405--1454.
\bibitem[Deryugina(2017)]{deryugina2017} Deryugina, T. (2017). ``The Fiscal Cost of Hurricanes: Disaster Aid versus Social Insurance.'' \textit{American Economic Journal: Economic Policy} 9(3): 168--198.
\bibitem[Dube et~al.(2020)]{dube2020} Dube, A., J. Jacobs, S. Naidu and S. Suri (2020). ``Monopsony in Online Labor Markets.'' \textit{American Economic Review: Insights} 2(1): 33--46.
\bibitem[Farber(2015)]{farber2015} Farber, H. S. (2015). ``Why You Can't Find a Taxi in the Rain and Other Labor Supply Lessons from Cab Drivers.'' \textit{Quarterly Journal of Economics} 130(4): 1975--2026.
\bibitem[Hsiang and Jina(2014)]{hsiang2014} Hsiang, S. M. and A. S. Jina (2014). ``The Causal Effect of Environmental Catastrophe on Long-Run Economic Growth: Evidence from 6,700 Cyclones.'' NBER Working Paper 20352.
\bibitem[Kahn(2005)]{kahn2005} Kahn, M. E. (2005). ``The Death Toll from Natural Disasters: The Role of Income, Geography, and Institutions.'' \textit{Review of Economics and Statistics} 87(2): 271--284.
\end{thebibliography}
\end{singlespace}

\end{document}
